\section{Proofs of the Security Statements}
\label{app:security-proofs}
These proofs use game $G_{\rm bind}$ and its exact winning event in
Section~\ref{sec:security-game}. S3 is an event-inclusion argument within
the stated experiment, not a UC composition theorem. They are separate from the existing Lean
artifact. They quantify over polynomially bounded adversaries and traces
under the stated encoding, key and infrastructure premises.

\subsection{Proof of Security Theorem~\ref{sec-thm:binding}}
\begin{proof}
Maintain, for every protected key $i$, the set $Q_i$ of message digests
submitted to its signing oracle, together with the typed envelope and all
original reference preimages for each submission. The simulator can
evaluate public signature verification and compare captured finite
inventories, so it can detect the first $\mathrm{BadBind}$ acceptance.
Let $(i,x^*,(b_\ell^*)_\ell,m^*,\sigma^*)$ be that acceptance.

If $m^*\notin Q_i$, verification under $pk_i$ succeeds on a message never
queried to that key's signing oracle. This is an EUF-CMA forgery, regardless
of whether a similar envelope or a signature under another key was seen.
For a single-key reduction choose a uniform protected-key index $j$ before
setup, embed the challenge public key there, generate all other keys and
simulate their signing queries locally. Forward only key $j$'s requests
to the challenge signing oracle. This is the original distribution because
the protected keys are independently generated and none is corrupted.
Stop at the first $\mathrm{BadBind}$ acceptance. If its digest is fresh
and its issuer is $j$, output
$(m^*,\sigma^*)$. The simulator has no need to predict the winning
verification attempt; in every other case it aborts. Its success probability
is the first win's fresh-digest case probability
divided by $N$, giving the term $N\epsilon_{\rm sig}$.

Otherwise $m^*\in Q_i$. Choose any recorded query with that digest and call
its envelope $x$. If $x^*\ne x$ as typed messages, injectivity gives
\[
 C(\mathrm{envelope},x^*)\ne C(\mathrm{envelope},x),
\]
while their hash outputs equal $m^*$. These two explicit preimages are
a collision for $H$. Different type, instance, role, epoch, predecessor
or inline payload fields are included in this case, not treated as
unprotected metadata.

If $x^*=x$, inventory names, kinds, positions and all commitments coincide.
If every referenced byte string also coincides with the chosen issuance
entry, that entry witnesses correspondence, contradicting
$\mathrm{BadBind}$. Therefore some referenced bytes $b_\ell^*\ne b_\ell$
differ. Acceptance recomputed the stored commitment from the captured
bytes, and issuance computed it from its retained original bytes. Thus
the distinct, explicitly available inputs
\[
 \begin{gathered}
 C(\mathrm{object},\mathrm{kind}_\ell,\mathrm{name}_\ell,b_\ell^*)
 \\
 C(\mathrm{object},\mathrm{kind}_\ell,\mathrm{name}_\ell,b_\ell)
 \end{gathered}
\]
have equal hash outputs. For a nested manifest, choose the first differing
manifest or descendant in the finite reference inventory; the same argument
applies. No inaccessible original preimage is invented by the reduction.

A collision reduction generates the signature keys itself and simulates the
whole experiment while retaining these inputs. It also stops at the first
$\mathrm{BadBind}$ acceptance, outputs the pair in the old-digest case
and aborts in the fresh-digest case. This single reduction covers all
typed envelope/object uses and all verification attempts. The fresh-digest
and old-digest cases exhaust the first win, giving the stated sum.
Repeated identical signing requests merely add matching log entries;
the theorem asserts existence of an earlier issuance, not uniqueness of
its signing event or signature representation.
\end{proof}

\subsection{Proof of Security Theorem~\ref{sec-thm:replay}}
\begin{proof}
Consider the linearized sequence of state-changing operations. Immediately
after the first successful operation with token $z$, atomic commitment has
placed $z$ in the persistent consumed-token set together with its effect.
Every later state retains $z$. A second success with $z$ would require both
membership (by persistence) and nonmembership (by its precondition), a
contradiction. A crash before commitment produces no committed effect; a
crash after commitment may leave the client uncertain but does not undo
token consumption. Receipt lookup after retry causes no second effect.

If two commits use predecessor $h$, the first successful one changes the
head to a value different from $h$ in the same operation. The no-head-reuse
premise prevents an ABA return to $h$. The other commit's comparison
therefore fails. Signature authenticity alone is irrelevant to this
ordering argument: even two genuinely issued requests cannot both pass
the same stale comparison.

After revocation, suspension or epoch replacement linearizes, current
status or epoch no longer satisfies the old authorization's grant guard.
At or beyond expiry, the strict time comparison also fails. Atomic grant
evaluation uses that current state, so an old authorization cannot produce
a later grant. A grant ordered earlier is already admitted; the argument
does not cancel in-flight responses. A fresh authorization may satisfy
the updated guards, which explains the theorem's explicit exception.

This proves at-most-once effect and current admission within the state
model. It does not derive state durability, correct clock observations,
linearizability or complete mediation from cryptographic assumptions.
\end{proof}

\subsection{Proof of Security Theorem~\ref{sec-thm:composition}}
\begin{proof}
Use S3's coupled execution with a persistent failure flag; stopping at its
first setting recovers the S1 game without changing the flag probability.
Let $V$ be the existence of an ordinary-RELEASE grant violating one of its
registered thresholds anywhere in the full execution. Let $F=\bigcup_j E_j^c$ for the coverage families
defined on this same adaptive experiment. On
$\mathrm{BadBind}^{\,c}\cap F^c$, S1's correspondence and the explicit
scientific premises give correctly bound, applicable upper endpoints for
every selected in-scope claim. The deterministic reducer gives
$\theta_{ctk}\le U_{ctk}\le\tau_{ctk}$ for ordinary RELEASE.
The current-state, atomic-grant and complete-mediation premises connect
that same candidate/context to each grant. A changed context cannot
inherit the earlier evidence without revalidation. Hence
\[
 V\subseteq\mathrm{BadBind}\cup F.
\]
Theorem~\ref{thm:statistics}'s conditional allocation argument gives
$\Pr(F)\le\alpha$ on the entire experiment. S1 gives the other term, so
subadditivity yields
\[
 \Pr(V)\le\Pr(\mathrm{BadBind})+\Pr(F)
       \le N\epsilon_{\rm sig}+\epsilon_H+\alpha.
\]
This reasoning never asserts
$\Pr(F\mid\mathrm{BadBind}^{\,c})\le\alpha$, which need not follow from
an unconditional coverage guarantee. For a bounded execution, unused
future allocations may be set to zero. No independence between signature,
hash and statistical events is used. If authentic measurement fabrication
or an unmediated output path is allowed, the event inclusion can fail;
the bound then does not apply rather than silently charging those failures
to the statistical budget.
\end{proof}
